\documentclass[12pt]{article}
\usepackage[utf8]{inputenc}
\usepackage[T1]{fontenc}
\usepackage[margin=1in]{geometry}
\usepackage{setspace}
\usepackage{xcolor}
\usepackage[colorlinks=true, urlcolor=blue!60!black, linkcolor=blue!60!black]{hyperref}
\setlength{\parindent}{2em}
\setlength{\parskip}{1.25em}
\renewcommand{\baselinestretch}{1.0}
\begin{document}
\begin{center}
{\LARGE \textit{The Prize Cases}, 67 U.S. (2 Black) 635 (1863)\par}
\vspace{0.8em}
{\large Case Brief --- Four-Page Edition\par}
\end{center}
\vspace{1em}

\section{Memory Jogger}

In a 5--4 wartime decision, the Court upheld President Lincoln's unilateral April 1861 naval blockade of Confederate ports, ruling that while the President cannot initiate or declare war, he is constitutionally bound to repel sudden invasions or armed rebellions by force without waiting for a congressional declaration of war (\href{https://supreme.justia.com/cases/federal/us/67/635/}{67 U.S. 635}).

\section{Detailed Case Facts}

After the Confederate bombardment of Fort Sumter on April 12, 1861, with Congress out of session, Lincoln invoked the Militia Act of 1795 and the Insurrection Act of 1807 and proclaimed a blockade of Southern ports on April 19, extended to Virginia and North Carolina on April 27. In special session, Congress passed the Act of July 13, 1861, authorizing the President to declare states in insurrection, and the Act of August 6, 1861, retroactively approving his prior proclamations.

Before July 13, Union vessels captured four merchant ships as prizes: the \textit{Amy Warwick}, a Richmond-owned brig carrying coffee from Rio de Janeiro, whose master was ignorant of the war; the \textit{Crenshaw}, a Richmond-owned schooner carrying tobacco to Liverpool; the \textit{Hiawatha}, a British barque carrying tobacco and cotton from Richmond; and the \textit{Brilliante}, a Mexican schooner that entered New Orleans after being warned off the Mississippi blockade.

\section{Procedural History}

The United States and the capturing crews filed prize libels in federal district courts in Massachusetts, New York, and Florida, which condemned the vessels and cargoes (restoring some loyal and neutral interests). The claimants appealed directly, arguing that no legal war existed before July 13, 1861, and that the President could not institute a blockade without a declaration of war. After twelve days of argument, the Court decided the cases on March 10, 1863.

\section{Judicial Votes}

Affirmed, \textbf{5--4}.
\begin{itemize}
\item \textbf{Majority:} Justice Grier, joined by Justices Wayne, Swayne, Miller, and Davis.
\item \textbf{Concurrences:} None.
\item \textbf{Dissent:} Justice Nelson, joined by Chief Justice Taney and Justices Catron and Clifford.
\end{itemize}

\section{Holding}

The condemnations were affirmed, except that tobacco bought before the war by New York citizens in \textit{The Crenshaw} was restored. Although the President cannot initiate or declare war, when invasion or armed rebellion of civil-war magnitude occurs he is authorized and bound to resist force by force, including by blockade, without awaiting a declaration of war. Whether an insurrection has reached that point is a political judgment to which courts defer, and property of persons domiciled in hostile territory is capturable as ``enemies' property'' regardless of personal loyalty. Grier's view that the August 1861 Act ``operated to perfectly cure the defect,'' if any, was an alternative ground.

\section{Analysis of Opinions}

\textbf{Grier (majority).} Drawing on Vattel and Lord Stowell, Grier held that war is a matter of fact: a civil war is ``never solemnly declared'' but exists when rebels hold territory, organize armies, and close the courts. Congress alone declares war, but the Executive Power, the Take Care Clause, and the Commander-in-Chief Clause, reinforced by the 1795 and 1807 Acts, require the President to meet war thrust upon the nation. Courts must defer to his judgment, and Queen Victoria's neutrality proclamation estopped neutrals from denying a war. In a territorial civil war bounded by ``lines of bayonets,'' the government holds both sovereign and belligerent rights, so commerce from hostile territory is capturable.

\textbf{Nelson (dissenting).} Nelson accepted that Congress recognized a civil war on July 13, 1861, but held all earlier captures void. He distinguished war ``in a material sense'' from war ``in a legal sense,'' which turns citizens into enemies, voids commerce, and authorizes confiscation, and which the Constitution commits to Congress. The 1795 and 1807 Acts let Lincoln wage a ``personal war'' against insurgents until Congress met, but not convert loyal citizens into enemies, and retroactive ratification was an impermissible ex post facto confiscation. The \textit{Hiawatha} was also entitled to a warning.

\section{Examples: Future Applications}

\begin{enumerate}
\item \textbf{Same-Side (Valid): Surprise Attack.} A foreign fleet strikes a Pacific naval base during a congressional recess, and the President orders the fleet destroyed and a maritime exclusion zone imposed. He did not initiate the war and is bound to resist force by force.
\item \textbf{Same-Side (Valid): Armed Insurrection.} Separatists seize federal forts and expel federal courts in two states; the President invokes the Insurrection Act and quarantines the territory. Courts defer to his assessment.
\item \textbf{Opposite-Side (Invalid): Preventive War.} The President invades a nation that might develop weapons in five years but has not attacked. Grier said the President ``has no power to initiate or declare a war.''
\item \textbf{Opposite-Side (Invalid): Domestic Seizures and Military Trials.} The President seizes steel mills or tries civilians militarily where courts are open, beyond the war zone, as in \textit{Ex parte Milligan} and \textit{Youngstown}.
\item \textbf{Fence-Sitter: Cyberattacks and Attacks on Allies.} A state-sponsored cyberattack disables a U.S. power grid, or missiles strike an ally's ships, and the President orders retaliatory airstrikes. Is this repelling a sudden attack or initiating war?
\end{enumerate}

\section{Critique}

\textbf{Originalist and textualist views.} Michael Ramsey and Saikrishna Prakash read the Declare War Clause as leaving the President only a narrow power to repel sudden attacks, making Grier's limit faithful to original meaning; John Yoo argues that Article II confers broad power to initiate hostilities unless Congress checks it. Others note that reliance on inherent Article II power obscured the Calling Forth Clause's placement of insurrection responses under congressional control.

\textbf{Doctrinal and separation-of-powers views.} Louis Fisher, David Barron, and Martin Lederman argue that administrations overread the decision as a blank check, ignoring that Lincoln acted defensively, with Congress out of session, and was promptly ratified. Deborah Pearlstein notes the Court still scrutinized individual captures and restored property in \textit{The Crenshaw}.

\section{Key Quotations}

\begin{quote}
``He has no power to initiate or declare a war either against a foreign nation or a domestic State. \ldots If a war be made by invasion of a foreign nation, the President is not only authorized but bound to resist force by force. He does not initiate the war, but is bound to accept the challenge without waiting for any special legislative authority.''
\end{quote}
\noindent --- Justice Grier, Opinion of the Court, 67 U.S. at 668.

\begin{quote}
``Whether the President, in fulfilling his duties as Commander-in-chief in suppressing an insurrection, has met with such armed hostile resistance and a civil war of such alarming proportions as will compel him to accord to them the character of belligerents is a question to be decided by him, and this Court must be governed by the decisions and acts of the political department of the Government to which this power was entrusted.''
\end{quote}
\noindent --- Justice Grier, Opinion of the Court, 67 U.S. at 670.

\begin{quote}
``The President was bound to meet it in the shape it presented itself, without waiting for Congress to baptize it with a name; and no name given to it by him or them could change the fact.''
\end{quote}
\noindent --- Justice Grier, Opinion of the Court, 67 U.S. at 669.

\end{document}
